% GENERATED FILE. Edit canonical lessons, then run npm run generate:books.
% canonical-source-hash: fnv1a64:7e0097b496da4fef
% canonical-lessons: KA-C436-hamali, KA-C436-trali, KA-C436-heru, KA-C436-kaisamanu, KA-C436-bhartiyagu

\chapter{Porters and Trolleys}
\label{ch:ka-porters-and-trolleys}
\hlchaptermodality{436}

\begin{chapteropening}
\textbf{By the end of this chapter:} \emph{I can say a porter, a trolley, to load, hand luggage and to fill up.}

Complete the last lesson of chapter 436: I can say a porter, a trolley, to load, hand luggage and to fill up.
\end{chapteropening}

\section[\texorpdfstring{\kn{ಹಮಾಲಿ} (hamāli)}{ಹಮಾಲಿ (hamāli)}]{\kn{ಹಮಾಲಿ} (hamāli) — a porter}

\label{lesson:KA-C436-hamali}



\begin{quote}
Before the new one: say the Kannada for to peep, then the Kannada for a fare.
\end{quote}

\subsection*{You'll want to know: \kn{ಹಮಾಲಿ}}
\textbf{\kn{ಹಮಾಲಿ}} — \emph{hamāli} — \textquotedblleft{}a porter\textquotedblright{}.

\begin{quote}
\kn{ಹಮಾಲಿ} \kn{ನನ್ನ} \kn{ಎರಡು} \kn{ಸೂಟ್ಕೇಸ್ಗಳನ್ನು} \kn{ಹೊತ್ತನು}.
\end{quote}

\begin{quote}
\emph{hamāli nanna eraḍu sūṭkēsgaḷannu hottanu.}
\end{quote}

\begin{quote}
The porter carried my two suitcases.
\end{quote}

Plural: \kn{ಹಮಾಲಿಗಳು}, \emph{hamāligaḷu}.

At stations porters wear a red shirt and a numbered badge; \kn{ಕೂಲಿ} also names them.

\begin{cousinweb}
\textbf{\kn{ಹಮಾಲಿ}} (\emph{hamāli}) is usually traced to Arabic \emph{ḥammāl}, \textquotedblleft{}a carrier, a porter\textquotedblright{}, which reached Kannada through Hindustani or Marathi.
\end{cousinweb}

\begin{grammarlens}[title={What you've built}]
One more word for getting around.
\end{grammarlens}

\subsection*{Your turn}
\begin{itemize}
\raggedright
  \item \emph{Say it:} \emph{hamāli}
  \item \emph{Say it:} \emph{hamāli}, once more
  \item \emph{Say it:} \emph{hamāli nanna eraḍu sūṭkēsgaḷannu hottanu.} — The porter carried my two suitcases.
  \item \emph{Recall:} say the Kannada for to peep, then the Kannada for a fare, then say \emph{hamāli} again
\end{itemize}

\begin{tcolorbox}[breakable,colback=teal!4,colframe=teal!35!black,title={Before you move on}]
What does \kn{ಹಮಾಲಿ} mean? (\textquotedblleft{}A porter\textquotedblright{}.) Say it once more.
\end{tcolorbox}
\section[\texorpdfstring{\kn{ಟ್ರಾಲಿ} (ṭrāli)}{ಟ್ರಾಲಿ (ṭrāli)}]{\kn{ಟ್ರಾಲಿ} (ṭrāli) — a trolley}

\label{lesson:KA-C436-trali}



\begin{quote}
Before the new one: say the Kannada for a fare, then the Kannada for a porter.
\end{quote}

\subsection*{You'll want to know: \kn{ಟ್ರಾಲಿ}}
\textbf{\kn{ಟ್ರಾಲಿ}} — \emph{ṭrāli} — \textquotedblleft{}a trolley\textquotedblright{}.

\begin{quote}
\kn{ವಿಮಾನ} \kn{ನಿಲ್ದಾಣದಲ್ಲಿ} \kn{ಟ್ರಾಲಿ} \kn{ಉಚಿತ}.
\end{quote}

\begin{quote}
\emph{vimāna nildāṇadalli ṭrāli ucita.}
\end{quote}

\begin{quote}
At the airport the trolley is free.
\end{quote}

A luggage trolley at a station, or a hand cart in a shop.

\begin{cousinweb}
\textbf{\kn{ಟ್ರಾಲಿ}} (\emph{ṭrāli}) is a loan from English \emph{trolley}, a small cart pushed by hand.
\end{cousinweb}

\begin{grammarlens}[title={What you've built}]
One more word for getting around.
\end{grammarlens}

\subsection*{Your turn}
\begin{itemize}
\raggedright
  \item \emph{Say it:} \emph{ṭrāli}
  \item \emph{Say it:} \emph{ṭrāli}, once more
  \item \emph{Say it:} \emph{vimāna nildāṇadalli ṭrāli ucita.} — At the airport the trolley is free.
  \item \emph{Recall:} say the Kannada for a fare, then the Kannada for a porter, then say \emph{ṭrāli} again
\end{itemize}

\begin{tcolorbox}[breakable,colback=teal!4,colframe=teal!35!black,title={Before you move on}]
What does \kn{ಟ್ರಾಲಿ} mean? (\textquotedblleft{}A trolley\textquotedblright{}.) Say it once more.
\end{tcolorbox}
\section[\texorpdfstring{\kn{ಹೇರು} (hēru)}{ಹೇರು (hēru)}]{\kn{ಹೇರು} (hēru) — to load; to impose}

\label{lesson:KA-C436-heru}



\begin{quote}
Before the new one: say the Kannada for a porter, then the Kannada for a trolley.
\end{quote}

\subsection*{You'll want to know: \kn{ಹೇರು}}
\textbf{\kn{ಹೇರು}} — \emph{hēru} — \textquotedblleft{}to load; to impose\textquotedblright{}.

\begin{quote}
\kn{ಅವರು} \kn{ವಾಹನದ} \kn{ಮೇಲೆ} \kn{ಸಾಮಾನು} \kn{ಹೇರಿದರು}.
\end{quote}

\begin{quote}
\emph{avaru vāhanada mēle sāmānu hēridaru.}
\end{quote}

\begin{quote}
They loaded the luggage on top of the vehicle.
\end{quote}

Past (I …): \kn{ಹೇರಿದೆ}, \emph{hēride}.

What is loaded takes -\kn{ಅನ್ನು} or no ending; figuratively, \kn{ತೆರಿಗೆ} \kn{ಹೇರು}, to impose a tax.

\begin{cousinweb}
\textbf{\kn{ಹೇರು}} (\emph{hēru}) is native Dravidian: DEDR 4446 lists Kannada \emph{pēṟu}, \textquotedblleft{}to lift up and put on, to load; a load\textquotedblright{}, the older form with the \emph{p-} that modern Kannada turns into \emph{h-}. The same article holds Telugu \emph{pērcu}, to pile up.
\end{cousinweb}

\begin{grammarlens}[title={What you've built}]
One more word for getting around.
\end{grammarlens}

\subsection*{Your turn}
\begin{itemize}
\raggedright
  \item \emph{Say it:} \emph{hēru}
  \item \emph{Say it:} \emph{hēru}, once more
  \item \emph{Say it:} \emph{avaru vāhanada mēle sāmānu hēridaru.} — They loaded the luggage on top of the vehicle.
  \item \emph{Recall:} say the Kannada for a porter, then the Kannada for a trolley, then say \emph{hēru} again
\end{itemize}

\begin{tcolorbox}[breakable,colback=teal!4,colframe=teal!35!black,title={Before you move on}]
What does \kn{ಹೇರು} mean? (\textquotedblleft{}To load; to impose\textquotedblright{}.) Say it once more.
\end{tcolorbox}
\section[\texorpdfstring{\kn{ಕೈಸಾಮಾನು} (kaisāmānu)}{ಕೈಸಾಮಾನು (kaisāmānu)}]{\kn{ಕೈಸಾಮಾನು} (kaisāmānu) — hand luggage}

\label{lesson:KA-C436-kaisamanu}



\begin{quote}
Before the new one: say the Kannada for a trolley, then the Kannada for to load.
\end{quote}

\subsection*{You'll want to know: \kn{ಕೈಸಾಮಾನು}}
\textbf{\kn{ಕೈಸಾಮಾನು}} — \emph{kaisāmānu} — \textquotedblleft{}hand luggage\textquotedblright{}.

\begin{quote}
\kn{ಕೈಸಾಮಾನು} \kn{ಏಳು} \kn{ಕಿಲೋಗಿಂತ} \kn{ಹೆಚ್ಚು} \kn{ಇರಬಾರದು}.
\end{quote}

\begin{quote}
\emph{kaisāmānu ēḷu kilōginta heccu irabāradu.}
\end{quote}

\begin{quote}
Hand luggage must not weigh more than seven kilos.
\end{quote}

\kn{ಕೈ}, hand, plus \kn{ಸಾಮಾನು}, luggage: the bag you keep with you on a plane or a bus.

\begin{cousinweb}
\textbf{\kn{ಕೈಸಾಮಾನು}} (\emph{kaisāmānu}) is a Kannada compound of native \emph{kai}, \textquotedblleft{}hand\textquotedblright{}, and \kn{ಸಾಮಾನು}, \textquotedblleft{}luggage, things\textquotedblright{}, which is usually traced to Persian \emph{sāmān}, \textquotedblleft{}goods, equipment\textquotedblright{}.
\end{cousinweb}

\begin{grammarlens}[title={What you've built}]
One more word for getting around.
\end{grammarlens}

\subsection*{Your turn}
\begin{itemize}
\raggedright
  \item \emph{Say it:} \emph{kaisāmānu}
  \item \emph{Say it:} \emph{kaisāmānu}, once more
  \item \emph{Say it:} \emph{kaisāmānu ēḷu kilōginta heccu irabāradu.} — Hand luggage must not weigh more than seven kilos.
  \item \emph{Recall:} say the Kannada for a trolley, then the Kannada for to load, then say \emph{kaisāmānu} again
\end{itemize}

\begin{tcolorbox}[breakable,colback=teal!4,colframe=teal!35!black,title={Before you move on}]
What does \kn{ಕೈಸಾಮಾನು} mean? (\textquotedblleft{}Hand luggage\textquotedblright{}.) Say it once more.
\end{tcolorbox}
\section[\texorpdfstring{\kn{ಭರ್ತಿಯಾಗು} (bhartiyāgu)}{ಭರ್ತಿಯಾಗು (bhartiyāgu)}]{\kn{ಭರ್ತಿಯಾಗು} (bhartiyāgu) — to fill up; to be full}

\label{lesson:KA-C436-bhartiyagu}



\begin{quote}
Before the new one: say the Kannada for to load, then the Kannada for hand luggage.
\end{quote}

\subsection*{You'll want to know: \kn{ಭರ್ತಿಯಾಗು}}
\textbf{\kn{ಭರ್ತಿಯಾಗು}} — \emph{bhartiyāgu} — \textquotedblleft{}to fill up; to be full\textquotedblright{}.

\begin{quote}
\kn{ಹಬ್ಬದ} \kn{ಕಾರಣ} \kn{ಎಲ್ಲಾ} \kn{ಬಸ್ಸುಗಳು} \kn{ಭರ್ತಿಯಾದವು}.
\end{quote}

\begin{quote}
\emph{habbada kāraṇa ellā bassugaḷu bhartiyādavu.}
\end{quote}

\begin{quote}
Because of the festival all the buses filled up.
\end{quote}

Past (it …): \kn{ಭರ್ತಿಯಾಯಿತು}, \emph{bhartiyāyitu}.

\kn{ಭರ್ತಿ} alone means full or filled; \kn{ಭರ್ತಿ} \kn{ಮಾಡು} is to fill in, as a form.

\begin{cousinweb}
\textbf{\kn{ಭರ್ತಿಯಾಗು}} (\emph{bhartiyāgu}) joins \emph{bharti}, \textquotedblleft{}full, filled\textquotedblright{}, usually traced to Hindustani \emph{bhartī}, \textquotedblleft{}filling, recruitment\textquotedblright{}, and the native verb \kn{ಆಗು}, to become.
\end{cousinweb}

\begin{grammarlens}[title={What you've built}]
One more word for getting around.
\end{grammarlens}

\subsection*{Your turn}
\begin{itemize}
\raggedright
  \item \emph{Say it:} \emph{bhartiyāgu}
  \item \emph{Say it:} \emph{bhartiyāgu}, once more
  \item \emph{Say it:} \emph{habbada kāraṇa ellā bassugaḷu bhartiyādavu.} — Because of the festival all the buses filled up.
  \item \emph{Recall:} say the Kannada for to load, then the Kannada for hand luggage, then say \emph{bhartiyāgu} again
\end{itemize}

\begin{tcolorbox}[breakable,colback=teal!4,colframe=teal!35!black,title={Before you move on}]
What does \kn{ಭರ್ತಿಯಾಗು} mean? (\textquotedblleft{}To fill up; to be full\textquotedblright{}.) Say it once more.
\end{tcolorbox}
